% Section content template for individual Educative lessons
% This file contains the actual content for a specific section
% Generated from Educative API section content

% Section: Key Insights and Tips: Designing a Parking Lot
% Section ID: 4610542210514944
% Section Slug: key-insights-and-tips-designing-a-parking-lot
% Generated: 2025-12-11T13:40:44.482526

You've now completed the parking lot case study, and this lesson builds on that foundation by addressing common design pitfalls and sharing effective interview techniques. To help reinforce your understanding, a short quiz and a selection of related case studies are included to further strengthen your grasp of object-oriented design concepts.

\subsection{Common mistakes}\label{r2gDRfoMdo1Ox8JatwR9Z}
\\
Avoiding the following common mistakes will help you develop a more efficient and scalable parking lot system:

\begin{itemize}
\item
\protect\phantomsection\label{feSt-UKhShPFq0bfcClTA}
\textbf{Ignoring multiple entry and exit points:} Failing to account for multiple entrances and exits in the initial design. Your \texttt{ParkingLot} class should manage collections of \texttt{Entrance} and \texttt{Exit} objects to accurately model the system.
\item
\protect\phantomsection\label{cG1gEkqr4hRwK4QszdQAx}
\textbf{Hardcoding parking spot allocation:} Do not create a flexible system for assigning spots to different vehicle types. Implement a strategy that allows a motorcycle to park in a compact or large spot if its designated spots are full.
\item
\protect\phantomsection\label{vHl-Jf9AbbGXsq-nxCmCO}
\textbf{Neglecting payment flexibility:} Designing a rigid payment system that only handles one method or point of payment. Using an abstract \texttt{Payment} class, the design should support various payment methods (cash, card) and locations (exit panel, parking agent).
\item
\protect\phantomsection\label{gDytjHlhTJqXoa2k46vg3}
\textbf{Designing a monolithic} \textbf{\texttt{ParkingLot}} \textbf{class:} Placing all logic and data directly within the main \texttt{ParkingLot} class makes it difficult to manage. Decompose the system into smaller, single-responsibility classes like \texttt{ParkingTicket}, \texttt{ParkingRate}, and \texttt{DisplayBoard}.
\item
\protect\phantomsection\label{KTtnYsF6qFi_mUn7KXjOO}
\textbf{Not tracking ticket status:} Lacking a mechanism to follow a parking ticket through its life cycle. Use a \texttt{TicketStatus} enumeration (e.g., issued, paid, canceled) to manage the state of each \texttt{ParkingTicket}.
\end{itemize}

\subsection{Interview tips}\label{GR1gPxd8ThX-062hRiPuy}
\\
Here are some essential tips to help you prepare for an interview where you're asked to design a parking lot:

\begin{center}
\renewcommand{\arraystretch}{1.5}

\begin{longtable}{|p{0.45\linewidth}|p{0.45\linewidth}|}
\hline
\textbf{Tip} & \textbf{Explanation} \\
\hline
\endfirsthead

\hline
\textbf{Tip} & \textbf{Explanation} \\
\hline
\endhead

\hline
\endfoot

\hline
\endlastfoot

\textbf{Clarify the types of vehicles and spots first.} &
Before writing any code, confirm the different kinds of vehicles and parking spots the system must support. This shows the interviewer you value requirement gathering and helps define the necessary class hierarchy (e.g., \texttt{Vehicle} $\rightarrow$ \texttt{Car}, \texttt{Truck}; \texttt{ParkingSpot} $\rightarrow$ \texttt{Compact}, \texttt{Large}). \\
\hline

\textbf{Justify your choice of class over enum.} &
Be prepared to defend your decision to use abstract classes for entities like \texttt{Vehicle} instead of enums. Explain that this adheres to the Open/Closed principle, making the system more scalable and maintainable if new vehicle types are introduced later. \\
\hline

\textbf{Define the \texttt{ParkingTicket} as a central class.} &
Emphasize the importance of the \texttt{ParkingTicket} class, explaining that it acts as a central data object linking the \texttt{Vehicle}, entry time, exit time, payment status, and total amount. \\
\hline

\textbf{Mention the need for multiple display boards.} &
A real-world system would have multiple display boards, likely one per floor or section. This demonstrates that you are thinking about the practical, physical implementation of the system. \\
\hline

\textbf{Address the ``full lot'' scenario.} &
Explain how your system will handle the case where the parking lot is full. Mention checks at the \texttt{Entrance} panels and updates to the \texttt{DisplayBoard} to show you've considered edge cases and system constraints. \\
\hline

\end{longtable}
\end{center}


\subsection{Test your understanding}\label{_Yr6Wm2koXiy540gGdiD6}
\\
Take a short quiz to evaluate your understanding of the parking lot design.

\subsection*{Designing a Parking Lot}
\vspace{0.3cm}
\noindent
\textbf{Question 1:} Why is using an abstract \texttt{Vehicle} class preferred over a \texttt{Vehicle} enum?
\\[0.2cm]
\begin{itemize}
 \item It is faster at runtime.
 \item It uses less memory.
 \item It is required for creating objects.
 \item \textbf{[CORRECT]} It supports the Open/Closed principle for better scalability.
\\[0.1cm]
\textit{Explanation:} 
An abstract class can be extended with new vehicle types (e.g., \texttt{Bus}) without modifying existing code, adhering to the Open/Closed principle.
\end{itemize}
\vspace{0.3cm}
\noindent
\textbf{Question 2:} What is the primary responsibility of the \texttt{Entrance} class?
\\[0.2cm]
\begin{itemize}
 \item To validate payment
 \item \textbf{[CORRECT]} To generate a parking ticket
\\[0.1cm]
\textit{Explanation:} 
The \texttt{Entrance} class is responsible for issuing a new \texttt{ParkingTicket} to a vehicle upon its arrival at the parking lot.
 \item To display free spots
 \item To calculate the parking rate
\end{itemize}
\vspace{0.3cm}
\noindent
\textbf{Question 3:} Which relationship does the ``Pay ticket'' use case have with the ``Scan ticket'' use case?
\\[0.2cm]
\begin{itemize}
 \item \textbf{[CORRECT]} Include
\\[0.1cm]
\textit{Explanation:} 
The diagram specifies an \textless{\textless{} include\textgreater{}\textgreater{}} relationship, as scanning the ticket to determine the fee is a necessary precondition for paying it.
 \item Extend
 \item Association
 \item Generalization
\end{itemize}
\vspace{0.3cm}
\noindent
\textbf{Question 4:} Which class is directly responsible for calculating the parking fee based on duration?
\\[0.2cm]
\begin{itemize}
 \item \texttt{Payment}
 \item \texttt{Exit}
 \item \textbf{[CORRECT]} \texttt{ParkingRate}
\\[0.1cm]
\textit{Explanation:} 
The \texttt{ParkingRate} class contains the logic to calculate the total charge based on the hours a vehicle has been parked.
 \item \texttt{ParkingTicket}
\end{itemize}
\vspace{0.3cm}
\noindent
\textbf{Question 5:} What kind of relationship exists between the \texttt{ParkingLot} class and the \texttt{ParkingSpot} class?
\\[0.2cm]
\begin{itemize}
 \item Inheritance
 \item Association
 \item Generalization
 \item \textbf{[CORRECT]} Composition
\\[0.1cm]
\textit{Explanation:} 
The \texttt{ParkingLot} comprises \texttt{ParkingSpot} objects; the spots cannot exist independently of the lot, indicating a strong composition relationship.
\end{itemize}

\subsection{Related case studies}\label{LQNaabpvsIjTl9sWj-xdA}
\\
Explore these related case studies to deepen your understanding of the design principles applied in the parking lot system:

\begin{itemize}
\item
\protect\phantomsection\label{HymjFV-WdhO9S87CBGInV}
\textbf{Designing a car rental system:} This case study also involves managing a fleet of vehicles, handling different vehicle types, and processing bookings and payments, reinforcing concepts of resource management and state tracking.
\item
\protect\phantomsection\label{Xu99cWLeTemrLuzb5xGux}
\textbf{Designing an elevator system:} This problem focuses on managing a limited set of resources (elevators) to efficiently serve user requests, similar to how a parking lot manages a limited number of spots.
\item
\protect\phantomsection\label{Eakqen9z8IwAIaQ9XNY3A}
\textbf{Designing a movie ticket booking system:} This involves managing seats (finite resources) in different screens and handling user transactions, sharing the core challenge of resource allocation and management with the parking lot design.
\item
\protect\phantomsection\label{taviRPRx7q2oRbpE8PHls}
\textbf{Designing a warehouse management system:} This requires tracking inventory in specific locations (bins/shelves), managing incoming and outgoing items, and optimizing placement, analogous to managing vehicles in parking spots.
\item
\protect\phantomsection\label{ng38ZBv59rAeXzxKaXWAH}
\textbf{Designing a toll booth system:} This case study focuses on tracking vehicles passing through specific points, calculating fees (often variable), and processing payments, sharing the transaction and vehicle tracking aspects.
\end{itemize}

% End of section content